% from ega1-5-fr.tex
\oldpage[I]{136}
\begin{corollary}[5.4.5]
\label{I.5.4.5-fr}
Soient $Z$ un $S$-schéma, $j:X\to Y$, $g:X\to Z$ deux $S$-morphismes. Si $j$
est une immersion fermée, il en est de même de
$(j,g)_S:X\to Y\times_S Z$.
\end{corollary}

La démonstration à partir de (\hyperref[I.5.4.4-fr]{5.4.4}) est la même que
celle de (\hyperref[I.5.3.14-fr]{5.3.14}) à partir de
(\hyperref[I.5.3.13-fr]{5.3.13}).

\begin{corollary}[5.4.6]
\label{I.5.4.6-fr}
Si $X$ est un $S$-schéma, toute $S$-section de $X$
(\hyperref[I.2.5.5-fr]{2.5.5}) est une immersion fermée.
\end{corollary}

Si $\varphi:X\to S$ est le morphisme structural, $\psi:S\to X$ une
$S$-section de $X$, il suffit d'appliquer
(\hyperref[I.5.4.5-fr]{5.4.5}) à $\varphi\circ\psi=1_S$.

\begin{corollary}[5.4.7]
\label{I.5.4.7-fr}
Soient $S$ un préschéma intègre, $s$ son point générique, $X$ un $S$-schéma.
Si deux $S$-sections $f$, $g$ de $X$ sont telles que $f(s)=g(s)$, alors
$f=g$.
\end{corollary}

En effet, si $x=f(s)=g(s)$, les homomorphismes $k(x)\to k(s)$ correspondant à
$f$ et $g$ sont nécessairement identiques. Si $h=(f,g)_S$, on en déduit
(\hyperref[I.5.3.17-fr]{5.3.17}) que $h(s)$ appartient à la diagonale
$Z=\Delta_X(X)$~; mais comme $S=\overline{\{s\}}$ et que $Z$ est fermée par
hypothèse, on a $h(S)\subset Z$. Il résulte alors de
(\hyperref[I.5.2.2-fr]{5.2.2}) que $h$ se factorise en
$S\to Z\to X\times_S X$, et on en conclut que $f=g$ par définition de la
diagonale.

\begin{remark}[5.4.8]
\label{I.5.4.8-fr}
Si l'on suppose réciproquement que la conclusion de
(\hyperref[I.5.4.3-fr]{5.4.3}) est vérifiée lorsque $f=1_Y$, on en conclut
que $Y$ est séparé au-dessus de $S$~; de même, si on suppose que la conclusion
de (\hyperref[I.5.4.5-fr]{5.4.5}) s'applique aux deux morphismes
\[
  Y\xrightarrow{\Delta_Y}Y\times_Z Y\xrightarrow{p_1}Y,
\]
on en tire que $\Delta_Y$ est une immersion fermée, donc que $Y$ est séparé
au-dessus de $Z$~; enfin, la validité de la conclusion de
(\hyperref[I.5.4.6-fr]{5.4.6}) pour la $Y$-section $\Delta_Y$ du
$Y$-préschéma $Y\times_S Y\to Y$, implique que $Y$ est séparé au-dessus de
$S$.
\end{remark}

\subsection{Critères de séparation.}
\label{subsection:I.5.5-fr}

\begin{proposition}[5.5.1]
\label{I.5.5.1-fr}
\begin{enumerate}
  \item[\rm(i)] Tout monomorphisme de préschémas (et en particulier toute
    immersion) est un morphisme séparé.
  \item[\rm(ii)] Le composé de deux morphismes séparés est séparé.
  \item[\rm(iii)] Si $f:X\to X'$, $g:Y\to Y'$ sont deux $S$-morphismes
    séparés, $f\times_S g$ est séparé.
  \item[\rm(iv)] Si $f:X\to Y$ est un $S$-morphisme séparé, le
    $S'$-morphisme $f_{(S')}$ est séparé pour toute extension $S'\to S$ du
    préschéma de base.
  \item[\rm(v)] Si le composé $g\circ f$ de deux morphismes est séparé, $f$
    est séparé.
  \item[\rm(vi)] Pour qu'un morphisme $f$ soit séparé, il faut et il suffit
    que $f_{\mathrm{red}}$ (\hyperref[I.5.1.5-fr]{5.1.5}) le soit.
\end{enumerate}
\end{proposition}

(i) résulte aussitôt de (\hyperref[I.5.3.8-fr]{5.3.8}). Si $f:X\to Y$,
$g:Y\to Z$ sont deux morphismes, le diagramme
\[
  \xymatrix{
    X\ar[rr]^{\Delta_{X|Z}}\ar[dr]_{\Delta_{X|Y}} & &
    X\times_Z X\\
    & X\times_Y X\ar[ur]_j
  }
  \tag{5.5.1.1}\label{I.5.5.1.1-fr}
\]
où $j$ désigne l'immersion canonique (\hyperref[I.5.3.10-fr]{5.3.10}) est
commutatif, comme on le vérifie aussitôt. Si $f$ et $g$ sont séparés,
$\Delta_{X|Y}$ est une immersion fermée par définition, et $j$ est une
immersion fermée en vertu de (\hyperref[I.5.4.2-fr]{5.4.2}), donc
$\Delta_{X|Z}$ est une immersion fermée par
(\hyperref[I.4.2.4-fr]{4.2.4}), ce qui

